\chapter{Generate the Store Application}

The store will render HTML for a browser, so it is not an API-only Rails
application. We will still use Rails' minimal preset: it retains controllers,
views, Active Record, and the asset pipeline while omitting JavaScript,
Hotwire, mail, jobs, deployment tooling, and most developer extras this
experiment does not need.

\section{Exercise 11A: Create the application account}

From the store jail's root console, create a dedicated group, user, and
application directory:

\begin{lstlisting}[language=sh]
pw groupadd store
pw useradd store -g store -m -s /bin/sh \
  -c "Store Rails application"
install -d -o store -g store -m 0755 \
  /usr/local/www/store
\end{lstlisting}

Confirm the ownership boundary before generating files:

\begin{lstlisting}[language=sh]
id store
stat -f '%Su:%Sg %Sp %N' /usr/local/www/store
\end{lstlisting}

The directory should belong to \texttt{store:store}.

\section{Exercise 11B: Generate a minimal full-stack application}

Run the generator as the application user:

\begin{lstlisting}[language=sh]
su -l store -c 'cd /usr/local/www/store && \
  rails new . \
    --minimal \
    --database=postgresql \
    --skip-bundle \
    --skip-git'
\end{lstlisting}

Unlike the catalog generator, this command does not use \texttt{--api}.
The store therefore keeps Action View and can render HTML. The
\texttt{--minimal} preset implies the skips for JavaScript, Hotwire, Docker,
Kamal, Solid Queue, and several other optional subsystems.

Verify that the generated files still belong to the service account:

\begin{lstlisting}[language=sh]
stat -f '%Su:%Sg %Sp %N' \
  /usr/local/www/store/Gemfile
\end{lstlisting}

\section{Exercise 11C: Satisfy the generated bundle}

The minimal full-stack application adds Puma and Propshaft. Its generated
development and test group also retains the Debug gem. Install all three as
prebuilt FreeBSD gems, then resolve the lockfile without downloading or
compiling anything:

\begin{lstlisting}[language=sh]
pkg install -y rubygem-puma rubygem-propshaft \
  rubygem-debug
su -l store -c \
  'cd /usr/local/www/store && bundle lock --local'
su -l store -c \
  'cd /usr/local/www/store && bundle check'
\end{lstlisting}

The final command should report that the bundle is satisfied.

\section{Checkpoint}

Stop when \texttt{bundle check} succeeds. The store application exists but is
not yet connected to PostgreSQL or the catalog service. Those integrations
will be added separately so each boundary can be tested on its own.
